\chapter{MARCO TEÓRICO}
\label{cap:teorico}

Este capítulo define los elementos teóricos que sostienen el diseño propuesto y
fija las \emph{fórmulas} de todas las métricas empleadas. El criterio es
explícito: una métrica no se reporta como un número aislado, sino junto a la
expresión que la define, al número de variables que la componen y a la
identificación de cuál de ellas domina el valor agregado.

\section{Segmentación semántica de nubes de puntos}

Los métodos supervisados actuales resuelven bien la tarea en dominios cerrados
---PTv3~\cite{ptv3} alcanza 77,5 mIoU y Sonata~\cite{sonata} 79,4 mIoU sobre
ScanNet--- pero comparten una restricción estructural: \emph{el conjunto de
clases se fija durante el entrenamiento}. Un elemento no previsto no se marca
como desconocido, sino que se fuerza hacia la clase conocida más parecida. En
patrimonio arquitectónico esa restricción es fatal, porque cada edificio
incorpora elementos singulares ---soluciones propias de su época, su taller o su
reconstrucción tras un sismo--- que no aparecen en ningún catálogo previo. Un
modelo de conjunto cerrado no puede señalar su propia ignorancia, y el
especialista pierde justamente la información que necesitaba: qué partes del
levantamiento requieren su atención.

\section{Descubrimiento de clases novedosas (NCD)}

Sea $P$ una nube de puntos y $\mathcal{C}=\mathcal{C}_k\cup\mathcal{C}_n$ el
conjunto de clases, con $\mathcal{C}_k\cap\mathcal{C}_n=\emptyset$. Durante el
entrenamiento se dispone de etiquetas únicamente para los puntos de
$\mathcal{C}_k$. El modelo debe producir una asignación
$\hat{y}:P\rightarrow\mathcal{C}$ correcta sobre $\mathcal{C}_k$ y consistente
sobre $\mathcal{C}_n$. Es, en esencia, la capacidad de detectar formas nuevas sin
haberlas visto etiquetadas.

\section{Prototipos e hipergrafo (HyperNCD)}

HyperNCD~\cite{hyperncd} representa cada punto mediante un descriptor compuesto
\begin{equation}
F_i=\left[f_{\mathrm{sem}}(x^i);\,f_{\mathrm{geo}}(x^i)\right],
\label{eq:hyperncd3}
\end{equation}
correspondiente a la Ec.~3 de~\cite{hyperncd}, donde $f_{\mathrm{sem}}$ proviene
de una red MinkowskiUNet-34C y $f_{\mathrm{geo}}$ se calcula analíticamente sobre
los $k$ vecinos más cercanos ($k=15$): linealidad, planaridad y \emph{scattering}.
Aguas abajo, la afinidad entre prototipos combina una componente geométrica y una
semántica,
\begin{equation}
S_b=\alpha\,S_g+(1-\alpha)\,S_s,\qquad \alpha\in[0,1],
\label{eq:afinidad}
\end{equation}
y cada prototipo conserva sus $M=8$ vecinos más afines al construir el
hipergrafo, sobre el que se propaga la información de alto orden que permite
agrupar las clases novedosas.

\section{Representación auto-supervisada y \emph{geometric shortcut}}

Sonata~\cite{sonata} provee un \emph{encoder} PTv3 preentrenado de forma
auto-supervisada que produce $f_{\mathrm{ssl}}(x^i)$ con pesos congelados y uso
exclusivo en inferencia. Su aportación conceptual es el diagnóstico del
\emph{geometric shortcut}: la representación colapsa sobre pistas geométricas
fácilmente accesibles ---normales, altura, curvatura local--- en lugar de
aprender semántica. La calidad de un descriptor congelado se mide por
\emph{linear probing}, entrenando únicamente un clasificador lineal $W$ sobre él,
\begin{equation}
\hat{y}_i=\arg\max\left(W^{\top}f_{\mathrm{ssl}}(x^i)+b\right),
\label{eq:probe}
\end{equation}
de modo que el desempeño resultante es atribuible a la representación y no al
clasificador. Esa propiedad es la que permite, más adelante, separar la
contribución de la representación de la del razonamiento por hipergrafo.
\section{Modelo híbrido propuesto}

La propuesta de este trabajo consiste en conservar íntegro el razonamiento por
hipergrafo de HyperNCD y sustituir su rama geométrica manual por las
representaciones aprendidas del \emph{encoder} congelado de Sonata, de modo que
\begin{equation}
F_i=\left[f_{\mathrm{sem}}(x^i);\,f_{\mathrm{ssl}}(x^i)\right].
\label{eq:hibrido}
\end{equation}
La mejora esperada es de métrica: un mIoU más alto sobre las clases novedosas se
traduce directamente en menos corrección manual por parte del ingeniero, que es
la condición para que la herramienta tenga valor práctico.

\section{Métricas de segmentación y su formulación}
\label{sec:metricas}

\subsection{IoU y mIoU}

Para una clase $c$, la intersección sobre unión es
\begin{equation}
\mathrm{IoU}_c=\frac{TP_c}{TP_c+FP_c+FN_c},
\label{eq:iou}
\end{equation}
donde $TP_c$, $FP_c$ y $FN_c$ son los puntos verdaderos positivos, falsos
positivos y falsos negativos de esa clase. La métrica agregada es
\begin{equation}
\mathrm{mIoU}(\mathcal{S})=\frac{1}{|\mathcal{S}|}\sum_{c\in\mathcal{S}}\mathrm{IoU}_c,
\qquad
\mathcal{S}\in\{\mathcal{C}_k,\;\mathcal{C}_n,\;\mathcal{C}\},
\label{eq:miou}
\end{equation}
que se reporta por separado como mIoU \emph{known}, mIoU \emph{novel} y mIoU
\emph{all}.

\subsection{Análisis de dominancia}

La Ec.~\eqref{eq:miou} es una media aritmética no ponderada, de modo que
\begin{equation}
\frac{\partial\,\mathrm{mIoU}}{\partial\,\mathrm{IoU}_c}=\frac{1}{|\mathcal{S}|}
\qquad \forall\, c\in\mathcal{S}:
\label{eq:dominancia}
\end{equation}
cada clase pesa igual con independencia de su número de puntos. Por ello un valor
agregado alto nunca se reporta solo: se acompaña del vector
$(\mathrm{IoU}_c)_{c\in\mathcal{S}}$, de su desviación estándar y de la
identificación de la clase dominante
\begin{equation}
c^{*}=\arg\max_{c\in\mathcal{S}}\bigl|\mathrm{IoU}_c-\mathrm{mIoU}(\mathcal{S})\bigr|,
\label{eq:dominante}
\end{equation}
para saber qué variable explica el resultado. Un mIoU de 0,99 no es informativo
por sí mismo; lo es solo cuando se conoce cuántas clases promedia y cuál de ellas
lo gobierna.

\subsection{Emparejamiento de clases novedosas}

Como las clases de $\mathcal{C}_n$ se descubren sin etiqueta, la correspondencia
entre agrupaciones predichas y clases reales se resuelve por asignación húngara:
\begin{equation}
\sigma^{*}=\arg\max_{\sigma\in\mathfrak{S}_{|\mathcal{C}_n|}}
\sum_{c\in\mathcal{C}_n}\mathrm{IoU}\bigl(\hat{y}=\sigma(c),\;y=c\bigr).
\label{eq:hungaro}
\end{equation}

\subsection{Métricas de calidad de agrupamiento}

Para evaluar la estructura descubierta con independencia del emparejamiento se
emplean el índice de Rand ajustado y la información mutua normalizada:
\begin{equation}
\mathrm{ARI}=
\frac{\sum_{ij}\binom{n_{ij}}{2}
-\Bigl[\sum_i\binom{a_i}{2}\sum_j\binom{b_j}{2}\Bigr]\big/\binom{n}{2}}
{\tfrac{1}{2}\Bigl[\sum_i\binom{a_i}{2}+\sum_j\binom{b_j}{2}\Bigr]
-\Bigl[\sum_i\binom{a_i}{2}\sum_j\binom{b_j}{2}\Bigr]\big/\binom{n}{2}},
\label{eq:ari}
\end{equation}
\begin{equation}
\mathrm{NMI}(U,V)=\frac{2\,I(U;V)}{H(U)+H(V)},
\label{eq:nmi}
\end{equation}
donde $n_{ij}$ es el número de puntos comunes entre la agrupación $i$ y la clase
$j$, $a_i$ y $b_j$ sus marginales, $I$ la información mutua y $H$ la entropía.

\section{Métricas propias del modelo híbrido}
\label{sec:metricas-hibrido}

Medir la métrica del híbrido no equivale a medir la de Sonata ni la de HyperNCD
por separado. Al combinar dos métodos, las ventajas y los modos de falla de cada
uno se transfieren al modelo resultante, y las métricas estándar pueden no
reflejar la mejora atribuible a la integración. Se definen por ello dos
indicadores adicionales y un protocolo de atribución.

\subsection{Ganancia absoluta}

\begin{equation}
\Delta\mathrm{mIoU}^{\mathrm{novel}}=
\mathrm{mIoU}^{\mathrm{novel}}_{\mathrm{hib}}-\mathrm{mIoU}^{\mathrm{novel}}_{\mathrm{base}}.
\label{eq:delta}
\end{equation}

\subsection{Brecha cerrada $G$ (métrica propuesta)}

La ganancia absoluta no es interpretable por sí sola, porque el techo alcanzable
depende de la sección. Se propone normalizarla por la brecha que separa a la
línea base de su propio desempeño sobre las clases que sí vio etiquetadas:
\begin{equation}
G=\frac{\mathrm{mIoU}^{\mathrm{novel}}_{\mathrm{hib}}-\mathrm{mIoU}^{\mathrm{novel}}_{\mathrm{base}}}
{\mathrm{mIoU}^{\mathrm{known}}_{\mathrm{base}}-\mathrm{mIoU}^{\mathrm{novel}}_{\mathrm{base}}}.
\label{eq:G}
\end{equation}
$G$ es adimensional y evalúa tres variables: $G=0$ indica que la sustitución no
aporta; $G=1$ indica que las clases novedosas alcanzan la calidad con que el
modelo segmenta las conocidas; $G<0$ indica degradación. Los tres términos se
reportan siempre por separado, de modo que un valor de $G$ nunca se lea sin
conocer cuál de ellos lo domina.

\subsection{Atribución por etapas}

El \emph{pipeline} propuesto es separable ---inferencia congelada de Sonata
primero, razonamiento por hipergrafo después---, por lo que la métrica se reporta
también en el punto intermedio: calidad del descriptor por \emph{linear probing}
(Ec.~\eqref{eq:probe}) antes de la formación de prototipos, y mIoU
(Ec.~\eqref{eq:miou}) después. Esa separación es la que permite atribuir una
mejora a la representación o al razonamiento, y no al sistema tratado como caja
negra. Si el \emph{pipeline} fuese unificado, habría que discutir primero la
adecuación misma de la métrica elegida.

\section{Módulos del proyecto}

El marco anterior se traduce en cuatro módulos, que son las columnas del diseño
desarrollado en el marco metodológico:

\begin{enumerate}
\item \textbf{Módulo~1 --- Preparación de datos.} Teselado en cinco secciones,
voxelización a 5\,cm y normalización de la nube TLS.
\item \textbf{Módulo~2 --- Rama semántica.} MinkowskiUNet-34C de HyperNCD, que
produce $f_{\mathrm{sem}}$.
\item \textbf{Módulo~3 --- \emph{Encoder} auto-supervisado.} PTv3 preentrenado de
Sonata, congelado y usado solo en inferencia, que produce $f_{\mathrm{ssl}}$.
\item \textbf{Módulo~4 --- Fusión y razonamiento.} Punto donde reside el aporte:
la Ec.~\eqref{eq:hibrido} reemplaza a la Ec.~\eqref{eq:hyperncd3} antes de la
formación de prototipos; el hipergrafo permanece intacto.
\end{enumerate}
